%! Author = Admin
%! Date = 17.04.2026
\section{Сбор датасета ограниченного семейства кривых}

Нейросетевая часть проекта не обучается напрямую на траекториях управления. Вместо этого формируется датасет состояний и целевых значений параметрической скорости $V^{\ast}$, которая максимальна, но при этом не нарушает условия устойчивого слежения. Генерация датасета реализована в \texttt{code/ml/dataset/build\_dataset.py}.

\subsection{Семейство допустимых кривых}

В датасет включаются только кривые, совместимые с базовым регулятором, то есть удовлетворяющие условию
\begin{equation}
    \|t(s)\| = \mathrm{const}.
\end{equation}
Используются три типа траекторий:
\begin{enumerate}
    \item прямые;
    \item окружности;
    \item спирали.
\end{enumerate}

Для генератора кривых задаются вероятности выбора типов $(0.1,\ 0.4,\ 0.5)$ соответственно. Такое распределение увеличивает долю траекторий с ненулевой кривизной, поскольку именно они ограничивают допустимую скорость движения и дают информативные примеры для обучения.

\subsection{Формирование обучающего примера}

Для каждой кривой выбираются точки $s \in [s_{\min}, s_{\max}]$, после чего в их окрестности формируется начальное состояние
\begin{equation}
    x_0 = x_{\mathrm{curve}}(s) + \Delta x,
\end{equation}
где $\Delta x$ --- малое случайное возмущение положения и скорости. Это необходимо, чтобы признаки в датасете не вырождались в нули и отражали реальные отклонения от траектории.

Для каждого состояния извлекается вектор признаков
\begin{equation}
    f =
    \begin{bmatrix}
        e_1 &
        e_2 &
        \dot{e}_2 &
        \|v\| &
        e_{\mathrm{head}} &
        \kappa_c &
        \kappa_{\max}
    \end{bmatrix}^{\top},
\end{equation}
где:
\begin{equation}
    e_{\mathrm{head}} = \arccos\left(
        \frac{v^{\top} t(s)}{\|v\| \, \|t(s)\|}
    \right),
\end{equation}
\begin{equation}
    \kappa_c(s) =
    \frac{\|t(s) \times t'(s)\|}{\|t(s)\|^3},
\end{equation}
\begin{equation}
    \kappa_{\max}(s) = \max_{\tau \in [s,\ s+\Delta s]} \kappa_c(\tau).
\end{equation}

Большинство признаков нормируются по ограничениям объекта:
\begin{equation}
    \tilde{e}_1 = \frac{e_1}{e_{1,\max}}, \qquad
    \tilde{e}_2 = \frac{e_2}{e_{2,\max}}, \qquad
    \tilde{\dot{e}}_2 = \frac{\dot{e}_2}{v_{\max}}, \qquad
    \widetilde{\|v\|} = \frac{\|v\|}{V_{\max}}.
\end{equation}

\subsection{Оракул выбора целевой скорости}

Целевая переменная $V_{\mathrm{opt}}$ для каждой записи ищется с помощью оракула. Оракул перебирает значения скорости в диапазоне
\begin{equation}
    V \in [V_{\min}, V_{\max}]
\end{equation}
с заданным шагом и выбирает наибольшее значение, при котором короткий rollout остается устойчивым:
\begin{equation}
    V_{\mathrm{opt}} = \max \left\{ V : \mathrm{Stable}(x_0,\mathrm{curve},V) = 1 \right\}.
\end{equation}

Устойчивость rollout определяется тремя критериями:
\begin{equation}
    \max_t |e_2(t)| \le e_{2,\max},
\end{equation}
\begin{equation}
    \max_t \|v(t)\| \le v_{\max},
\end{equation}
\begin{equation}
    x(t) \ \text{не содержит NaN и Inf}.
\end{equation}

В конфигурации проекта для оракула по умолчанию используются параметры
\begin{equation}
    \Delta t_{\mathrm{oracle}} = 0.01, \qquad \kappa_{\mathrm{oracle}} = 100,
\end{equation}
а горизонт rollout автоматически выбирается из длины участка между соседними sample-точками, чтобы дрон успевал пройти достаточно информативный отрезок траектории.

\subsection{Структура датасета}

Сохраняемая запись CSV содержит поля
\begin{equation}
    \{e_1, e_2, \dot{e}_2, \|v\|, e_{\mathrm{head}}, \kappa_c, \kappa_{\max}, s, \|t\|, V_{\mathrm{opt}}\}.
\end{equation}
Из них первые семь признаков используются как вход нейросети, а $V_{\mathrm{opt}}$ является целевой переменной. Поля $s$ и $\|t\|$ сохраняются для диагностики и последующего анализа распределения примеров по кривым.

\subsection{Предварительная обработка и разбиение}

Перед обучением выполняются следующие шаги подготовки:
\begin{enumerate}
    \item валидация кривой на соответствие ограничению $\|t(s)\|=\mathrm{const}$;
    \item генерация возмущенных состояний вокруг кривой;
    \item извлечение и нормализация признаков;
    \item поиск $V_{\mathrm{opt}}$ с помощью оракула;
    \item отбрасывание неустойчивых или некорректных rollout;
    \item сохранение итоговой таблицы в \texttt{dataset.csv}.
\end{enumerate}

При обучении нейросети датасет далее делится на обучающую и валидационную части в отношении $0.8/0.2$. Такое разбиение реализовано в \texttt{train\_model.py}.
